\section{Introduction}
\label{sec:introduction}

The projection form of the Kadison--Kaplansky conjecture asserts that
the reduced group \(C^*\)-algebra of a torsion-free discrete group
contains no projections other than zero and the identity.
For a discrete group \(H\), this algebra is the operator-norm closure
of its left regular group algebra on \(\ell^2(H)\), and its canonical
trace is \(\tau_H(x)=\langle x\delta_1,\delta_1\rangle\).
The conjecture concerns scalar projections in this completion.
We prove the following result.

\begin{theorem}
\label{thm:main}
There exist a finitely generated torsion-free discrete group \(\Gproj\)
and a projection \(e\in C_r^*(\Gproj)\) such that
\[
 0<\tau_{\Gproj}(e)<\frac12.
\]
In particular, \(e\) is neither zero nor the identity.
\end{theorem}

Theorem~\ref{thm:main} resolves the projection form of the
Kadison--Kaplansky conjecture negatively.
The group is given by an infinite graphical presentation followed
by a free product with a finitely generated free group.
Every parameter is finite, although no useful bound on its size
is sought. The choices of labeling tables and edge voltages are
specified by finite distributions and proved simultaneous
existence statements.

\subsection*{Context}

Kadison traces his free-group projection problem to Kaplansky's question
about idempotents in simple unital \(C^*\)-algebras
\cite[pp.~222--224]{Kadison2008}.
Pimsner and Voiculescu proved the free-group case by computing the
\(K\)-theory of reduced crossed products by free groups
\cite[p.~131]{PimsnerVoiculescu1982}.
A central route to projectionlessness is integrality of the canonical
trace on \(K_0(C_r^*(H))\), where the trace on matrices is the
\emph{unnormalized} sum of the diagonal traces. For torsion-free groups,
the trace is integral on the image of reduced assembly with trivial
coefficients \cite[Theorem~0.3]{Luck2002},
\cite[Proposition~10.3]{BGW2016}. Surjectivity therefore gives scalar
projectionlessness. The survey \cite[Section~4.5.3]{GJV2019} explains the
distinct scalar and stable formulations.

This implication covers substantial classes of groups. Higson and
Kasparov proved Baum--Connes with coefficients for groups acting properly
by affine isometries on Hilbert space, including countable amenable groups
\cite[Theorem~9.1 and Corollary~9.2]{HK2001}. Mineyev and Yu established
the coefficient-free hyperbolic case and deduced the projection conjecture
for torsion-free subgroups of hyperbolic groups
\cite[Theorems~20--21]{MY2002}; Lafforgue later established the hyperbolic
case with arbitrary coefficients \cite[Theorem~0.4]{Lafforgue2012}.
The group constructed here has a different geometry: its presentation
encodes an infinite tower through embedded nilpotent covers.

For the group \(\Gproj\) in Theorem~\ref{thm:main}, the complex group
ring \(\mathbb C[\Gproj]\) has no scalar idempotents other than zero and
one \cite[Corollary~1.2(ii)]{OpenAIBass2026}. Thus the projection \(e\)
lies in the reduced completion but not in the group ring. The companion
also gives integer trace range on algebraic \(K_0(\mathbb C[\Gproj])\)
and on its image in \(K_0(C_r^*(\Gproj))\) under the comparison map
induced by inclusion \cite[Corollary~1.2(i) and Remark~10.1]{OpenAIBass2026}.
This does not assert integrality on the whole completed \(K_0\), where
\([e]\) has nonintegral trace.

Expanders also underlie the obstructions of Higson, Lafforgue,
and Skandalis for groupoids, coarse spaces, and group actions
with coefficients \cite{HLS2002}.
Their group-action statements concern assembly with coefficients,
whereas the trace obstruction here concerns trivial coefficients.

The later analysis of expander index obstructions by Willett and Yu
\cite{WY2012} gives another relevant spectral and coarse-geometric
precedent. Its group consequences likewise use coefficients.

Austin constructed rational group-ring operators with irrational
von Neumann kernel dimension \cite{Austin2013}.
Such a dimension alone does not imply that the kernel projection
lies in the reduced \(C^*\)-algebra.
Here a \emph{uniform separated spectral band} for a finite matrix
group-algebra operator in its regular representation gives
reduced-algebra membership by continuous functional calculus.
Scalar compression is applied after this matrix projection
has been obtained.

\subsection*{The obstruction and the proof strategy}

The first task is to obtain a matrix projection whose unnormalized trace
is strictly between zero and \(1/2\). A spectral projection in the group
von Neumann algebra is insufficient. We construct a finite self-adjoint
matrix \(T\) over the group algebra of a torsion-free group \(\Gbase\)
whose regular spectrum has a nonempty band near \(+1\), separated from
all remaining spectrum. Continuous functional calculus then places its
upper projection \(p_0\) in \(M_N(C_r^*(\Gbase))\). The difficult part
is to retain that gap while making the unnormalized trace small.

Begin with finite sets \(A_i,B_i\), indexed by \(i\in\mathbb Z\),
and bipartite incidences between \(A_i\) and \(B_k\) for \(|i-k|\le1\).
Their normalized operators preserve constant vectors and contract the
orthogonal complements by a common positive amount. Two versions of the
same tower serve different purposes: \(Y^0\) provides expansion and
\(Y^\#\) provides recovery of long paths from words in a finite alphabet.
They agree after deleting a uniformly small fraction of vertices.
The survivor \(Y\) is the base of a voltage cover \(P\): its vertices
are pairs \((v,f)\), with \(f\) in a torsion-free nilpotent group, and an
edge changes the second coordinate by its assigned group element.
The voltages give large girth and will also realize the required operators.

The analytic model is indexed by a central Fourier variable \(t\).
Twisted-convolution projections with effective parameters \(t\) and
\(t-\delta\), for a small positive \(\delta\), give a nonzero fiber
projection for \(0<t<\delta\). Normalized combinations of the constant
vectors on at most two adjacent layers produce an eigenvalue \(+1\),
while expansion separates all the other positive spectrum.
As \(t\) approaches either endpoint, these supporting layers move toward
an end of the tower and their fiber trace decreases. A finite magnetic
representation is essential to obtain small trace at both endpoints.
The result is a summable trace estimate even after the geometric weights
needed for intersections of copies have been included.

One edge voltage and one matrix weight per edge approximate this model
in the full regular norm. The same finite sampling construction enforces
the girth required for the group presentation. Long-word recovery then
embeds components of \(P\) as indexed copies in the Cayley graph of
\(\Gbase\). A common-root order confines every earlier intersection
with a copy to one finite tree ball. Removing those intersections leaves
disjoint domains; the remaining Cayley edges form a forest. The trace
estimate controls the loss in the upper spaces caused by the removed
balls, and a forest norm estimate controls the unused edges. Together
they preserve the separated band for the finite matrix \(T\).
For the trace estimate, we restrict a vector to every indexed copy and
project to its local upper space. This map is equivariant, although the
assignment of domains need not be. Disjointness and the same weighted
loss estimate bound its retained and removed parts. Its injectivity on
the global upper space then bounds the unnormalized trace by the total
trace of the local upper spaces.

The final passage to a scalar projection is a separate free-product
argument. Scale and comparison in reduced free products were studied by
Dykema and R{\o}rdam \cite{DR1998,DR2000,DRII2000}.
Recent work on selfless reduced free products and comparison by trace
\cite{FKOP2026,Robert2025} also underlies a conditional argument of
Sauers \cite{Sauers2026} that turns a nonintegral matrix projection over
\(C_r^*(H)\) into a scalar projection over \(C_r^*(H*\mathbb Z)\).
Such a reduction still requires an actual nonintegral matrix projection.
Here Proposition~\ref{scalar:compression} gives a direct quantitative
estimate at trace below \(1/2\): adjoining \(N\) free generators
compresses \(p_0\) to a scalar projection of the same trace. This gives
\(\Gproj=\Gbase*\mathbb F_N\), preserving finite generation and
torsion-freeness. The same small trace also places \([p_0]\) outside
the coefficient-free reduced assembly image for \(\Gbase\).

\subsection*{Organization}

Section~\ref{sec:scalar} reduces Theorem~\ref{thm:main} to the small
unnormalized matrix-trace target stated in Section~\ref{sec:setup}.
Sections~\ref{tow:section}--\ref{prune:section} construct the auxiliary
graph \(Y^0\) with expanding incidences, the gated graph \(Y^\#\)
with path recovery, and their common pruned induced subgraph \(Y\).

Section~\ref{geo:section} proves the embedding, overlap, forest, and
torsion conclusions under a girth hypothesis on the voltage cover
\(P\). Section~\ref{nil:section} constructs the deck group and the
small-trace fiber projections. Section~\ref{sam:section} identifies
the model's upper band, then chooses voltages and matrix weights
satisfying both the girth hypothesis and the regular operator
approximation.

Section~\ref{spec:section} passes the upper spectral band from the
model to the cover and then to the Cayley operator, and bounds its
unnormalized trace using all indexed copies. Section~\ref{sec:conclusion}
applies scalar compression to complete the proof.
